\section{Related work}\label{sec:related}

\paragraph{Zero-shot voice cloning}
Codec language models \citep{chen2025valle}, factorised diffusion \citep{ju2024naturalspeech3}, flow matching
\citep{chen2025f5tts} and large-language-model backbones with speech tokens \citep{du2024cosyvoice,du2024cosyvoice2,
du2025cosyvoice3,hu2026qwen3tts,liao2024fishspeech,casanova2024xtts} clone a voice from a few seconds of reference
audio. Several open models list Arabic among their languages: \vox, a tokenizer-free 2B-parameter model trained on a
multilingual corpus and released under Apache-2.0 \citep{zhou2025voxcpm,zhou2026voxcpm2,openbmb2026voxcpm2card};
Chatterbox Multilingual (MIT) \citep{resemble2025chatterbox}; and OmniVoice, which covers hundreds of languages under
a non-commercial licence \citep{zhu2026omnivoice}. Their reports evaluate Arabic, if at all, on MSA test sets.

\paragraph{Egyptian and dialectal Arabic speech synthesis}
Doan et al.\ \citep{doan2024arabicdialects} fine-tune XTTS for Arabic dialects and observe that the fine-tunes
produce dialectal speech while their WER rises; they read the rise as a cost. Habibi \citep{chen2026habibi} trains
unified dialectal models and reports WER against two recognisers next to a dialect MOS from listeners; its Egyptian
results include a system with lower WER and a lower dialect MOS than a commercial system, which the authors do not
discuss. NileTTS \citep{khamis2026niletts} builds a 38-hour synthetic Egyptian corpus from two voices and evaluates
with Whisper WER and speaker similarity only. Community fine-tunes of Chatterbox and OmniVoice for Egyptian
\citep{namaa2026egyptiantts,oddadmix2026lahgtna,aly2026voicetut} report no dialect measurement. None of these works
compares systems on intelligibility and dialect at once, or asks which one a selection rule would choose.

\paragraph{Evaluating dialect and accent}
Recognition error rates are known to be unreliable for under-represented accents. Zhong et al.\ \citep{zhong2025accent}
state that WER has ``significant limitations'' for assessing such accents and recommend against WER, CER and UTMOS
for accent generation; their argument rests on English with an Edinburgh accent and on an assumed ranking of systems,
not on listener ratings of intelligibility. Their follow-up \citep{zhong2026segmental} shows that segment-level
listener annotations can rank systems differently from similarity ratings. Classifier-based evaluation of dialect
and accent is proposed by Voxlect \citep{feng2025voxlect}, which reports agreement with native ratings for
Mandarin dialects, and by Vox-Profile \citep{feng2025voxprofile} for accent preservation in voice conversion. For
Arabic, dialect identification has been benchmarked in the MGB challenges \citep{ali2017mgb3,ali2019mgb5}; we use
two public classifiers trained on different data \citep{abdullah2025tamyizcard,abdullah2025adivc,
ibrahimamin2025hubertadi}. Automatic MOS predictors are known to transfer poorly to new languages and domains
\citep{huang2022voicemos,cooper2023voicemos,kunesova2024zeroshotmos}, so we report DNSMOS \citep{reddy2021dnsmos,
reddy2022dnsmosp835} only as a secondary signal. We are not aware of prior work showing, for any dialect, that the
system with the lowest recognition error is the one judged least dialectal by classifiers and listeners alike, so
that selection by intelligibility picks the wrong model.

\paragraph{Watermarking and consent}
AudioSeal \citep{sanroman2024audioseal} and WavMark \citep{chen2023wavmark} embed detectable marks in generated
speech; AudioSeal is localised and fast, and is what \BA uses. Hutiri et al.\ \citep{hutiri2024notmyvoice} catalogue
harms of speech generators, including impersonation and non-consensual cloning, which motivate the release choices in
\cref{sec:release}.
